\section{Temporal order, dynamic prediction and intervention regimes in R}
\label{exec14:r-temporal}

A treasury analyst may ask which funding action today improves tomorrow's
contribution, while a monitoring analyst asks which variables forecast tomorrow's
shortfall. The same rectangular file can support both questions, but the required
information differs. A row might be an independent institution with variables
measured at three waves, a daily observation from one institution, or an
observation under a known experimental regime. A tier orders variables within a
row; a lag relates rows; a regime describes how a mechanism changes. Giving these
objects distinct types is the first requirement for a useful R interface.

The inherited pcalg pathway covers an important part of this territory. Its
GIES interface estimates an interventional essential graph, using intervention
targets and a row-to-target index. An empty target represents observation without
intervention. The returned representative DAG is one member of the estimated
class. It must not silently become the uniquely established mechanism. Likewise,
an undirected essential-graph edge represents alternative directions across
DAGs: a plotting convention that renders it with two arrowheads does not turn
it into a latent-confounding PAG edge. FCI-JCI instead augments system variables
with context variables and explicitly selects context assumptions, including
exogeneity and, when invoked, absence of latent confounding between contexts and
the system. A month-end dummy does not establish these assumptions merely by
being supplied to the API. These are distinct statistical specifications, even
though both can combine data from multiple regimes. \citep{exec13_pcalg,exec13_latentr_primary}

\subsection{Ordering variables is different from modelling a series}

The inspected causalDisco 1.2.1 interface exposes TPC, TFCI and TGES with tiered
background knowledge. Its examples of temporally ordered variables describe a
useful cohort-wave application: measurements at an earlier wave cannot be caused
by measurements at a later wave under the declared clock. That restriction does
not supply stationarity, equal time intervals, independent observations or a
valid conditional-independence test for dependent rows. The selected TPC/TFCI
implementation wraps pcalg operations and converts their results through caugi.
The method name therefore describes a proposed workflow entry point, not a
complete uncertainty contract. \citep{exec14_rtime_causaldisco}

Two details require qualification before conversion. First, the pinned TFCI
help describes a PAG stored with class UNKNOWN because native PAG support is
unavailable, whereas the runner in that same revision requests class PAG from
the caugi constructor. The TPC help describes a PDAG while its runner requests
UNKNOWN. These are observed source/documentation disagreements; without the
matching dependency and a runtime inspection we cannot assert the actual returned
class or whether conversion succeeds. The proposed adapter should preserve the
original endpoint matrix and labels, record the effective constructor and
dependency version, and refuse an unsupported conversion. Reading an UNKNOWN
tag as an ordinary undirected graph would discard scientifically relevant
information. \citep{exec14_rtime_causaldisco}

Second, the tier-aware test wrapper returns zero without calling the underlying
test when a proposed conditioning variable is later than both endpoints. This
sentinel prevents that forbidden conditioning set from removing an edge. It is
not an observed p-value of zero. A diagnostic screen that reports every returned
number as a statistical test would misstate the evidence. The interface must
distinguish a test executed on data from a restriction imposed by declared
knowledge, just as it distinguishes an oriented arrow from a fitted coefficient.
TFCI's source places an arrowhead at the later endpoint of a cross-tier edge;
it does not thereby prove the earlier endpoint is an unconfounded causal parent.
The remaining endpoint and the equivalence-class semantics must survive export.
\citep{exec14_rtime_causaldisco}

The separate tpc 1.0 package is a useful comparator. Its API specifies a partial
node order, forbidden edges and global or tier-specific context variables,
returning a pcalg-style pcAlgo object. Its source forces particular context
adjacencies and removes others; supplied forbidden edges can override those
context constraints with a reported conflict. Thus labelling a business variable
as context can insert substantive structure rather than simply add a column to
the data. A reviewer should see the effective constraints after precedence rules,
their business rationale and which arrows follow only from those constraints.
This tpc package and causalDisco's TPC pathway have different package identities;
their names alone establish neither identical implementations nor identical
scientific behavior. \citep{exec14_rtime_tpc}

For cohort applications, we would first use such tools as comparators with
independent subjects and an explicit measurement order. Repeated accounts sharing
calendar shocks require an additional sampling argument. For example, suppose
100 banks have residual \(R_{it}=C_t\) on each of 20 dates, with independent
mean-zero, variance-one \(C_t\) across dates. The grand mean is the mean of the
20 calendar shocks. Its variance is \(1/20\), whereas an independent-row formula
using 2,000 rows gives \(1/2000\). The standard errors differ by a factor of ten.
This exact special case does not assign a universal effective sample size to
panels. It demonstrates why adding a time ordering to graph search cannot repair
a wrong replication unit. \citep{exec14_rtime_exact}

\subsection{Dynamic Bayesian networks and admissible rows}

dbnR 0.8.0 provides a different route: Gaussian dynamic Bayesian-network learning
with an explicit number of unrolled slices. Its size argument includes the
present slice, so size two corresponds to one historical lag. The structure
wrapper permits a previously folded data set through f\_dt and forwards it to
the selected search method; the reviewed wrapper does not by itself establish
the guarantees of each downstream search engine. This is a promising prediction
and lag-construction baseline. A learned DBN or a Gaussian conditional forecast
still requires causal assumptions before it supports an intervention query.
\citep{exec14_rtime_dbnr}

The data preparation source supplies a concrete safeguard worth reusing.
Ordinary fold\_dt shifts rows and removes incomplete leading rows.
filter\_same\_cycle compares the identifiers across slices, and
filtered\_fold\_dt combines folding with that filter. The integration obligation
is to select the appropriate pathway, supply the correct identifiers, and verify
the rows it produces. The selected functions do not establish chronological
sorting or equal physical intervals from timestamps. Equal bank identifiers
are necessary but insufficient for a one-day lag when days are missing.
\citep{exec14_rtime_dbnr}

Consider the deliberately small table of bank A values \(1,2,3\) on days
\(1,2,3\), followed by bank B values \(101,102,103\) on the same three days.
Positional one-lag folding produces five pairs, including the invalid transition
from bank A's last observation to bank B's first. Requiring equal bank identity
and a one-day increment retains four pairs. Conversely, two observations from
bank A on days one and three pass an identity-only filter but do not form a
one-day transition. These are exact constructions of required adapter behavior,
not reports of running dbnR or demonstrations that its documented ID filter
fails. \citep{exec14_rtime_exact}

The same source contains reduce\_freq, which averages batches whose size is
the ceiling of the requested interval divided by the current interval, with
optional identifier grouping. A physical-calendar adapter must consequently
record the achieved spacing and aggregation convention, especially when the
ratio is noninteger or timestamps are irregular. Averaging is not merely a
computational reduction: it changes the variables and potentially the estimand.
\citep{exec14_rtime_dbnr}

An exact daily example makes the change visible. Let the action \(A_t\) be
independent Rademacher variables, taking values minus one and one with equal
probability, and let the outcome be \(Y_t=A_{t-1}\). At the daily scale the
structural response is a one-day lag. With nonoverlapping two-day averages,
\[
 \bar A_k=\frac{A_{2k-1}+A_{2k}}2,
 \qquad
 \bar Y_k=\frac{A_{2k-2}+A_{2k-1}}2.
\]
Independence gives \(\operatorname{Var}(\bar A_k)=1/2\), while both
\(\operatorname{Cov}(\bar A_k,\bar Y_k)\) and
\(\operatorname{Cov}(\bar A_{k-1},\bar Y_k)\) equal \(1/4\).
Because successive action-block averages are independent, the population linear
projection of \(\bar Y_k\) on current and previous action-block averages has
coefficients one half and one half. A daily lag thus appears at both current
and previous block positions. \citep{exec14_rtime_exact}

Even the meaning of an intervention on a block average is incomplete without
its within-block timing. Increasing the first day's action by two changes the
current action average by one and the current outcome average by one. Increasing
the second day's action by two makes the same change to the current action
average, but changes the next outcome average by one instead. These perturbations
have identical aggregate action size and different response timing. A service
cannot answer an aggregate-action query with a unique response path until it
specifies which intervention is intended. The source data, aggregation map,
information clock and action policy all belong in the request identity.
\citep{exec14_rtime_exact}

\subsection{Invariant prediction: the exact scope of the guarantee}

seqICP 1.1 targets a response's invariant parents rather than returning a complete
temporal graph. Its linear API distinguishes an iid model from an autoregressive
model; lag order may be supplied or selected up to a maximum by AIC. It also
provides a nonlinear interface with an additive iid-noise model and a regression
function intended to approximate the target relation. These alternatives require
different justification. The manuals explicitly warn that reducing max.parents
below the predictor count sacrifices the stated confidence-interval coverage,
and that the default early stop when the accepted intersection is empty loses
coverage for coefficient values different from zero. The source also distinguishes
an empty intersection from modelReject, which records rejection of every tested
model. Both can display an empty parent set, yet communicate different evidence.
\citep{exec14_rtime_seqicp}

The genuine sequential-invariance original explains why these distinctions
matter. Its Definition 2.1 assumes a fixed linear Gaussian target relation with
iid innovations, and its causal interpretation additionally requires structural
independence conditions. Suppose \(S^*\) is the invariant true parent set and
\(\mathcal A\) is the family of accepted candidate sets. The estimator
\(\widehat S=\bigcap_{S\in\mathcal A}S\), with the declared empty-family
convention, satisfies
\[
 \Pr(\widehat S\subseteq S^*)\geq 1-\alpha
\]
when the relevant tests have the required level. If the true set is accepted,
its presence in the intersection ensures containment; erroneous inclusion of a
nonparent therefore requires rejecting that true set. This is a guarantee against
false included parents, not a guarantee of recovering every parent, selecting the
whole graph, or covering an intervention value with a numerical interval.
Failure to discover an edge is not a certificate that its effect is zero.
\citep{exec14_rtime_pfister}

Section 6 of that original adds the past \(p\) lags of all variables to each
candidate regression and targets instantaneous parents under a fixed linear
Gaussian higher-order Markov relation. The level statement is given as the
Monte Carlo resampling count \(B\) tends to infinity. It does not justify calling
the package's finite default \(B=100\) an exact finite-simulation guarantee.
The paper explicitly discusses lag selection and post-selection difficulties;
one proposed alternative considers multiple lag choices with a multiplicity
adjustment. AIC is a useful operational option, but its mere availability does
not remove that inferential issue. Moreover, Remark 3.2 labels the residual
permutation extension to nonlinear or non-Gaussian settings heuristic. We must
not borrow the linear fixed-lag theorem for every supplied nonlinear regressor.
\citep{exec14_rtime_pfister,exec14_rtime_seqicp}

Invariance also changes the usual conversation about stationarity. A stable
target mechanism can coexist with changes elsewhere in the system; those changes
may help rule out incorrect parent sets. Full-system stationarity is not a
universal desirable prerequisite for this route. A bank policy revision that
directly changes the target mechanism, however, can violate the very invariance
being tested. The specification should name which mechanisms may change, which
remain fixed, how blocks or smoothing comparisons are chosen, and whether
calendar shocks enter the innovation or observed predictors. An invariant
coefficient fit under an unjustified partition cannot establish that these
declarations are true. \citep{exec14_rtime_pfister}

\subsection{A funding counterexample and integration choices}

Temporal precedence and excellent prediction remain insufficient for the funding
question. Let \(C_t\) and \(E_t\) be mutually independent Rademacher sequences,
with the calendar driver \(C_t\) unobserved. In both of two worlds the recorded
action is \(A_t=C_t\). In the first world,
\(Y_{t+1}=2C_t+E_{t+1}\); in the second,
\(Y_{t+1}=2A_t+E_{t+1}\). Their complete stationary observed processes agree:
each daily pair \((A_t,Y_{t+1})\) takes
\((-1,-3),(-1,-1),(1,1),(1,3)\) with probability one quarter. Conditional
forecasting and temporal order cannot distinguish the worlds. Yet replacing
the action equation with \(A_t=1\) versus \(A_t=-1\) changes expected next-day
outcome by zero in the first world and four in the second. If the outcome is
stipulated to be thousands of dollars of funding contribution and that action
change costs one thousand dollars, the net consequences are minus one and plus
three thousand dollars. These values illustrate a missing-information decision,
not an estimate of bank return. \citep{exec14_rtime_exact}

The example is intentionally elementary. The practical implication is a typed
pipeline: separate graph candidates, predictive responses, identified effect
sets and monetary scenarios. Preserve interventions as mechanism changes rather
than row labels; retain lag direction, contemporaneous endpoints, measurement
frequency, within-period timing, regime membership and shared-calendar structure.
Return unknown when the requested effect lacks identification, and preserve
model rejection, unsupported conversion and insufficient test precision as
different reasons. A graph with plausible arrows should not bypass these checks.
\citep{exec14_rtime_contract}

Our proposed reuse choice is conditional integration of causalDisco for tiered
graph exploration, dbnR for checked lag/forecast baselines, and seqICP for a
declared target-parent invariance service. The separate tpc and inherited pcalg
paths remain reference comparators. No adapter is enabled by this comparison.
Direct adoption waits for dependency rights, pinned execution and semantic
qualification; porting is not selected because it would add a second
implementation before the original contract is understood. Declared licenses
are GPL-2 for the inspected causalDisco release, GPL-3 for dbnR and seqICP,
and GPL version three or later for tpc. Their dependency and distribution
obligations need a separate review; a process boundary is not assumed to settle
those questions. \citep{exec14_rtime_causaldisco,exec14_rtime_dbnr,exec14_rtime_seqicp,exec14_rtime_tpc,exec14_rtime_contract}

A useful first prototype can use public data for measurement conventions and
generated mechanisms for known truth. Predeclare tests for separate-series
boundaries, irregular calendars, aggregation timing, graph-endpoint round trips,
forced context constraints, all-rejected models, lag selection and finite
simulation settings. Compare a prediction-only baseline with correctly declared
and deliberately misspecified causal requests, keeping refusals and failures in
the denominator. The original R engines must actually run before we claim
runtime behavior or measured accuracy. Further screening should cover additional
nonlinear and mixed-frequency methods and the tsGFCI/Tetrad interface family;
the bounded repository-name search did not resolve that family's R route.
None of these local omissions blocks developing the remaining proposal areas.
\citep{exec14_rtime_contract}
\label{exec14:r-temporal-end}
